% !PRINTED 681
\weberpage{681}
%<<CONT
cess of the number $x$ over the greatest integer contained in $x$,
the following results:

If $\xi_1, \xi_2, \ldots, \xi_{\nu-1}$ is the system of exponents of a unit, and
$m$ an integer, then

\weq{6}{(m\,\xi_1), (m\,\xi_2), \ldots, (m\,\xi_{\nu-1})}

is likewise the system of exponents of a unit.

Now the quantities $(m\,\xi_i)$, when they are not zero, are
positive proper fractions, and by 1) it must therefore happen, if the
sequence of integers is continued far enough, that for two
different values $m'$, $m''$ of $m$ the sequence of numbers (6)
coincides. Since $m'\xi_i$, $m''\xi_i$ accordingly have the same
excess over an integer, $(m' - m'')\,\xi_i$ is itself an
integer, and there is therefore an integral rational
number $m$, at most equal to the number of systems of exponents
of reduced units, with the property that

\[
m\,\xi_1, m\,\xi_2, \ldots, m\,\xi_{\nu-1}
\]

become integral rational numbers.

This number $m$ may change when the unit $\varepsilon$ is
changed. But since all possible values of this number lie below a
finite limit, there is a finite number in which
all these values of $m$ are contained and which can itself be
taken for $m$. There is therefore a certain positive
integer $m$ that can be taken as the denominator of all those rational
 numbers which can occur as exponents of
 units. Theorem 6. is thereby proved.

\begin{center}
\S.\,187.\\[4pt]
\gesperrt{Fundamental Systems of Units.}
\end{center}

In place of the system of independent units $\varepsilon_i$ we now
 choose another, whose conjugate logarithms shall be

\weq{1}{\lambda_{i,1}, \lambda_{i,2}, \ldots, \lambda_{i,\nu} \qquad i = 1, 2, \ldots, \nu-1}

By \S.\,186, (5) we have

\weq{2}{\begin{aligned}
&\lambda_{i,k} = \xi_{1,i}\, l_{1,k} + \xi_{2,i}\, l_{2,k} + \cdots + \xi_{\nu-1,i}\, l_{\nu-1,k} \\
&i = 1, 2, \ldots, \nu-1; \quad k = 1, 2, \ldots, \nu,
\end{aligned}}

where $\xi_{1,i}, \xi_{2,i}, \ldots, \xi_{\nu-1,i}$ denote the exponents of one of the new units
 with respect to the system $\varepsilon_1, \varepsilon_2, \ldots, \varepsilon_{\nu-1}$, which
we may assume to be rational numbers with the common denominator $m$.
